\documentclass{article}
\usepackage[final]{nips_2017}
\usepackage[utf8]{inputenc} % allow utf-8 input
\usepackage[T1]{fontenc}    % use 8-bit T1 fonts
\usepackage{hyperref}       % hyperlinks
\usepackage{url}            % simple URL typesetting
\usepackage{booktabs}       % professional-quality tables
\usepackage{amsfonts}       % blackboard math symbols
\usepackage{nicefrac}       % compact symbols for 1/2, etc.
\usepackage{microtype}      % microtypography
\usepackage{graphicx}
\title{Formatting instructions for CS230-Fall 2019}

\author{
  John K.~Doe\thanks{Use footnote for providing further
    information about author (webpage, alternative
    address)---\emph{not} for acknowledging funding agencies.} \\
  Department of Computer Science\\
  Stanford University\\
  \texttt{***@stanford.edu} \\
  %% examples of more authors
  %% \And
  %% Coauthor \\
  %% Affiliation \\
  %% Address \\
  %% \texttt{email} \\
  %% \AND
  %% Coauthor \\
  %% Affiliation \\
  %% Address \\
  %% \texttt{email} \\
  %% \And
  %% Coauthor \\
  %% Affiliation \\
  %% Address \\
  %% \texttt{email} \\
  %% \And
  %% Coauthor \\
  %% Affiliation \\
  %% Address \\
  %% \texttt{email} \\
}

\begin{document}
% \nipsfinalcopy is no longer used

\begin{center}
\includegraphics[width=3cm, height=0.7cm]{CS230}
\end{center}

\maketitle

\begin{abstract}
The abstract should consist of 1 paragraph describing the motivation for your paper and a high-level explanation of the methodology you used/results obtained.
\end{abstract}

\section{Introduction}	
Explain the problem and why it is important. Discuss your motivation for pursuing this
problem. Give some background if necessary. Clearly state what the input and output
is. Be very explicit: “The input to our algorithm is an {image, amplitude, patient age,
rainfall measurements, grayscale video, etc.}. We then use a {SVM, neural network, linear
regression, etc.} to output a predicted {age, stock price, cancer type, music genre, etc.}.”
This is very important since different teams have different inputs/outputs spanning different
application domains. Being explicit about this makes it easier for readers. If you are using
your project for multiple classes, add a paragraph explaining which components of the
project were used for each class.

\section{Related work}
You should find existing papers, group them into categories based on their approaches,
and discuss their strengths and weaknesses, as well as how they are similar to and differ
from your work. In your opinion, which approaches were clever/good? What is the stateof-the-art?
Do most people perform the task by hand? You should aim to have at least
5 references in the related work. Include previous attempts by others at your problem,
previous technical methods, or previous learning algorithms. Google Scholar is very useful
for this: https://scholar.google.com/ (you can click “cite” and it generates MLA, APA,
BibTeX, etc.)

\section{Dataset and Features}
Describe your dataset: how many training/validation/test examples do you have? Is there
any preprocessing you did? What about normalization or data augmentation? What is the
resolution of your images? How is your time-series data discretized? Include a citation on
where you obtained your dataset from. Depending on available space, show some examples
from your dataset. You should also talk about the features you used. If you extracted
features using Fourier transforms, word2vec, PCA,
ICA, etc. make sure to talk about it. Try to include examples of your data in the report
(e.g. include an image, show a waveform, etc.).



\section{ Methods }
Describe your learning algorithms, proposed algorithm(s), or theoretical proof(s). Make
sure to include relevant mathematical notation. For example, you can include the loss function you are using. It is okay to use formulas from the lectures (online or in-class). For each algorithm, give a short description 
of how it works. Again, we are looking for your understanding of how these deep
learning algorithms work. Although the teaching staff probably know the algorithms, future
readers may not (reports will be posted on the class website). Additionally, if you are
using a niche or cutting-edge algorithm (anything else not covered in the class), you may want to explain your algorithm using 1/2
paragraphs. Note: Theory/algorithms projects may have an appendix showing extended
proofs (see Appendix section below).

\section{Experiments/Results/Discussion}
You should also give details about what (hyper)parameters you chose (e.g. why did you
use X learning rate for gradient descent, what was your mini-batch size and why) and how
you chose them. What your primary metrics are: accuracy, precision,
AUC, etc. Provide equations for the metrics if necessary. For results, you want to have a
mixture of tables and plots. If you are solving a classification problem, you should include a
confusion matrix or AUC/AUPRC curves. Include performance metrics such as precision,
recall, and accuracy. For regression problems, state the average error. You should have
both quantitative and qualitative results. To reiterate, you must have both quantitative
and qualitative results! If it applies: include visualizations of results, heatmaps,
examples of where your algorithm failed and a discussion of why certain algorithms failed
or succeeded. In addition, explain whether you think you have overfit to your training set
and what, if anything, you did to mitigate that. Make sure to discuss the figures/tables in
your main text throughout this section. Your plots should include legends, axis labels, and
have font sizes that are legible when printed.

\section{Conclusion/Future Work }
Summarize your report and reiterate key points. Which algorithms were the highestperforming?
Why do you think that some algorithms worked better than others? For
future work, if you had more time, more team members, or more computational resources,
what would you explore?

\section{Contributions}
The contributions section is not included in the 5 page limit. This section should describe
what each team member worked on and contributed to the project.

\section*{References}
This section should include citations for: (1) Any papers mentioned in the related work
section. (2) Papers describing algorithms that you used which were not covered in class.
(3) Code or libraries you downloaded and used. This includes libraries such as scikit-learn, Tensorflow, Pytorch, Keras etc. Acceptable formats include: MLA, APA, IEEE. If you
do not use one of these formats, each reference entry must include the following (preferably
in this order): author(s), title, conference/journal, publisher, year. If you are using TeX,
you can use any bibliography format which includes the items mentioned above. We are excluding
the references section from the page limit to encourage students to perform a thorough
literature review/related work section without being space-penalized if they include more
references. Any choice of citation style is acceptable
as long as you are consistent. 

\medskip
\small
[1] Alexander, J.A.\ \& Mozer, M.C.\ (1995) Template-based algorithms
for connectionist rule extraction. In G.\ Tesauro, D.S.\ Touretzky and
T.K.\ Leen (eds.), {\it Advances in Neural Information Processing
  Systems 7}, pp.\ 609--616. Cambridge, MA: MIT Press.

[2] Bower, J.M.\ \& Beeman, D.\ (1995) {\it The Book of GENESIS:
  Exploring Realistic Neural Models with the GEneral NEural SImulation
  System.}  New York: TELOS/Springer--Verlag.

[3] Hasselmo, M.E., Schnell, E.\ \& Barkai, E.\ (1995) Dynamics of
learning and recall at excitatory recurrent synapses and cholinergic
modulation in rat hippocampal region CA3. {\it Journal of
  Neuroscience} {\bf 15}(7):5249-5262.

\end{document}